\documentclass[a4paper]{arrowhead}

\usepackage[yyyymmdd]{datetime}
\usepackage{etoolbox}
\usepackage[utf8]{inputenc}
\usepackage[T1]{fontenc}
\usepackage{multirow}
\usepackage{longtable}
\usepackage{tikz}
\usetikzlibrary{positioning,shapes.multipart,arrows.meta}

\renewcommand{\dateseparator}{-}
\setlength{\emergencystretch}{3em}
\lstset{language=HTTP,breakatwhitespace=false,columns=fullflexible}

%% Special references
\newcommand{\mref}[2][]{{\textcolor{ArrowheadPurple}{\hyperref[sec:model:#2]{\ifstrempty{#1}{#2}{#1}}}}}
\newcommand{\oref}[1]{{\textcolor{ArrowheadBlue}{\hyperref[sec:operations:#1]{#1}}}}
\newcommand{\pdef}[1]{{\textcolor{ArrowheadGrey}{#1\label{sec:model:primitives:#1}\label{sec:model:primitives:#1s}\label{sec:model:primitives:#1es}}}}
\newcommand{\pref}[1]{{\textcolor{ArrowheadGrey}{\hyperref[sec:model:primitives:#1]{#1}}}}

% as template_sd.tex; the input and output arguments are passed already formatted (\mref or \pref)
\newrobustcmd\osubsection[3]{
  \addtocounter{subsection}{1}
  \addcontentsline{toc}{subsection}{\protect\numberline{\thesubsection}operation \textcolor{ArrowheadBlue}{#1}}
  \subsection*{
    \thesubsection\quad
    operation
    \textcolor{ArrowheadBlue}{#1}
    (\notblank{#2}{#2}{})
    \notblank{#3}{: #3}{}
  }
  \label{sec:operations:#1}
}
\newrobustcmd\msubsection[2]{
  \addtocounter{subsection}{1}
  \addcontentsline{toc}{subsection}{\protect\numberline{\thesubsection}#1 \textcolor{ArrowheadPurple}{#2}}
  \subsection*{\thesubsection\quad#1 \textcolor{ArrowheadPurple}{#2}}
  \label{sec:model:#2} \label{sec:model:#2s} \label{sec:model:#2es}
}
%%

\usepackage{xurl}
\begin{document}
\sloppy

%% Arrowhead Document Properties
\ArrowheadTitle{globalServiceDiscovery}
\ArrowheadServiceID{globalServiceDiscovery}
\ArrowheadType{Service Description}
\ArrowheadTypeShort{SD}
\ArrowheadVersion{5.3.0}
\ArrowheadDate{2026-11-09}
\ArrowheadAuthor{Jerker Delsing}
\ArrowheadStatus{REVIEW DRAFT}
\ArrowheadContact{jerker.delsing@ltu.se}
\ArrowheadFooter{\href{www.arrowhead.eu}{www.arrowhead.eu}}
\ArrowheadSetup
%%

%% Front Page
\begin{center}
  \vspace*{1cm}
  \huge{\arrowtitle}

  \vspace*{0.2cm}
  \LARGE{\arrowtype}
  \vspace*{1cm}

  \Large{Service ID: \textit{"\arrowid"}}

  \vspace*{\fill}

  % Front Page Image

  \vspace*{1cm}
  \vspace*{\fill}

  % Front Page Abstract
  \begin{abstract}
    This document describes the \texttt{global\allowbreak{}Service\allowbreak{}Discovery} service, version 5.2.2, produced by the Gatekeeper support system of Arrowhead (documentation edition 5.3.0).
  \end{abstract}

  \vspace*{1cm}
\end{center}

\newpage
%%

%% Table of Contents
\tableofcontents
\newpage
%%

\section{Documentation edition and implementation scope}
This is the 5.3.0 documentation edition, dated 2026-11-09. It carries forward the system and service contracts from the retained 5.2.2 documentation and identifies the current primary implementation sources. The former \path|documentation/v5.2.2| folder is retained unchanged.

Document version 5.3.0 is distinct from executable, dependency, registered system/service and wire-contract versions. Values still defined by the source (including 5.2.2 artefact versions, 1.0.0 service versions and \path|SysMonitor-AMQP-IDD-5.2.2| payload identifiers) are preserved. This documentation update does not change runtime behaviour or establish a new implementation release. Historical verification statements describe the tests recorded in the inherited baseline; this revision validates document compilation and references.

\section{Overview}
\label{sec:overview}

This document describes the \texttt{global\allowbreak{}Service\allowbreak{}Discovery} service, which provides discovery of authorized providers of a service in the neighbour clouds, through the inter-cloud relay.
It is produced by the Gatekeeper system \cite{gk-sysd} in service version 5.2.2.

The rest of this document is organized as follows.
In Section \ref{sec:operations}, we describe the abstract message operations provided by the service.
In Section \ref{sec:model}, we end the document by presenting the data types used by the mentioned operations.
The concrete interface is specified in the HTTP/TLS/JSON IDD \cite{gk-globalServiceDiscovery-http-idd}.

\subsection{Significant Prior Art}

Successor of the Arrowhead 4.6 Gatekeeper global service discovery (GSD); compatible with its paths (\texttt{/\allowbreak{}init\_\allowbreak{}gsd}).

\subsection{How This Service Is Meant to Be Used}

Dynamic\allowbreak{}Service\allowbreak{}Orchestration calls \texttt{query} during inter-cloud orchestration with the service requirement and the consumer. The Gatekeeper asks the Gatekeepers of the target clouds over the relay; each answers only with the providers its Consumer\allowbreak{}Authorization grants to the consumer of this cloud.

\subsection{Important Delimitations and Dependencies}
\label{sec:delimitations}

Only configured peer clouds (\texttt{FEDERATION\_\allowbreak{}PEER\_\allowbreak{}ORGANIZATIONS}) are asked; the answering cloud's identity must match. The profile is \texttt{legacy-\allowbreak{}mtls}: TLS per relay connection, not end-to-end encryption.

\newpage

\section{Service Interface}
\label{sec:operations}

This section describes the interfaces to the \texttt{global\allowbreak{}Service\allowbreak{}Discovery} service.
In particular, each subsection names an abstract operation, an input type and an output type, in that order.
The input type is named inside parentheses, while the output type is preceded by a colon.
Input and output types are only denoted when accepted or returned, respectively, by the interface in question.

All abstract data types named in this section are defined in Section \ref{sec:model}.

\begin{figure}[ht!]
  \centering
  \begin{tikzpicture}[font=\small]
    \node[draw=ArrowheadBlue, very thick, rectangle split, rectangle split parts=2, align=center,
          rectangle split part fill={ArrowheadBlue!12,white}, text width=9cm, inner sep=6pt] (b)
      {\guillemotleft{}Service\guillemotright{}\\\textbf{globalServiceDiscovery}
        \nodepart{second} \texttt{query}\\ \texttt{multi-\allowbreak{}query}};
  \end{tikzpicture}
  \caption{SysML-style block of the \texttt{global\allowbreak{}Service\allowbreak{}Discovery} service and its operations.}
  \label{fig:SD}
\end{figure}

The following interface operations are available.

\osubsection{query}{\mref[Global\allowbreak{}Service\allowbreak{}Discovery\allowbreak{}Request]{GlobalServiceDiscoveryRequest}}{\mref[Global\allowbreak{}Service\allowbreak{}Discovery\allowbreak{}Result]{GlobalServiceDiscoveryResult}}

Asks the target clouds for authorized providers of the required service.

\osubsection{multi-query}{\mref[Global\allowbreak{}Service\allowbreak{}Discovery\allowbreak{}Request]{GlobalServiceDiscoveryRequest}}{\mref[Global\allowbreak{}Service\allowbreak{}Discovery\allowbreak{}Result]{GlobalServiceDiscoveryResult}}

Same as \texttt{query} (v4 multi-GSD path).

\newpage

\section{Information Model}
\label{sec:model}

Here, all data objects that can be part of the \texttt{global\allowbreak{}Service\allowbreak{}Discovery} service are listed in alphabetic order.
Note that each subsection, which describes one type of object, begins with the \textit{struct} keyword, which is used to denote a collection of named fields, each with its own data type.
As a complement to the explicitly defined types in this section, there is also a list of implicit primitive types in Section \ref{sec:model:primitives}, which are used to represent things like hashes and identifiers.

\msubsection{struct}{GlobalServiceDiscoveryRequest}




\noindent\begin{tabularx}{\textwidth}{| >{\raggedright\arraybackslash}p{3.9cm} | >{\raggedright\arraybackslash\small}p{5.3cm} | X |} \hline
\rowcolor{gray!33} Field & Type & Description \\ \hline
service\allowbreak{}Requirement & \pref{Map}$<$\allowbreak{}\pref{String}, \pref{Object}$>$ & Orchestration service requirement; \texttt{serviceDefinition} is required, \texttt{operations} are checked (alias \texttt{requestedService}). \\ \hline
target\allowbreak{}Clouds & \pref{List}$<$\allowbreak{}\pref{String}$>$ & Clouds to ask (\texttt{\textless{}Cloud\textgreater{}} or \texttt{\textless{}Cloud\textgreater{}\textbar{}\allowbreak{}\textless{}Organization\textgreater{}}); default all configured peer clouds. \\ \hline
consumer\allowbreak{}System & \pref{String} & Consumer on whose behalf the orchestrator asks; default the requester. \\ \hline
\end{tabularx}



\msubsection{struct}{GlobalServiceDiscoveryResult}




\noindent\begin{tabularx}{\textwidth}{| >{\raggedright\arraybackslash}p{3.9cm} | >{\raggedright\arraybackslash\small}p{5.3cm} | X |} \hline
\rowcolor{gray!33} Field & Type & Description \\ \hline
results & \pref{List}$<$\allowbreak{}\mref{InterCloudProvider}$>$ & Authorized providers of all asked clouds. \\ \hline
unsuccessful\allowbreak{}Requests & \pref{Number} & Number of clouds that did not answer correctly. \\ \hline
queried\allowbreak{}Clouds & \pref{Number} & Number of clouds asked. \\ \hline
\end{tabularx}



\msubsection{struct}{InterCloudProvider}




\noindent\begin{tabularx}{\textwidth}{| >{\raggedright\arraybackslash}p{3.9cm} | >{\raggedright\arraybackslash\small}p{5.3cm} | X |} \hline
\rowcolor{gray!33} Field & Type & Description \\ \hline
cloud & \pref{Name} & Provider cloud, \texttt{\textless{}Cloud\textgreater{}\textbar{}\allowbreak{}\textless{}Organization\textgreater{}}. \\ \hline
provider\allowbreak{}Name & \pref{Name} & Provider system. \\ \hline
service\allowbreak{}Definition & \pref{Name} & Service definition. \\ \hline
version & \pref{Version} & Service version. \\ \hline
metadata & \pref{Metadata} & Instance metadata. \\ \hline
operations & \pref{List}$<$\allowbreak{}\pref{String}$>$ & Offered operations. \\ \hline
\end{tabularx}



\subsection{Primitives}
\label{sec:model:primitives}

Types and structures mentioned throughout this document that are assumed to be available to implementations of this service.
The concrete interpretations of each of these types and structures must be provided by any IDD document claiming to implement this service.

\noindent\begin{tabularx}{\textwidth}{| p{4.6cm} | X |} \hline
\rowcolor{gray!33} Type & Description \\ \hline
\pdef{List} & An \textit{array} of a known number of items, each having type A. \\ \hline
\pdef{Map} & A set of key-value pairs. \\ \hline
\pdef{Metadata} & A set of key-value pairs with arbitrary, possibly nested values. \\ \hline
\pdef{Name} & A string identifier that is intended to be both human and machine-readable. \\ \hline
\pdef{Number} & An integer number. \\ \hline
\pdef{Object} & Any structured value. \\ \hline
\pdef{String} & A chain of characters. \\ \hline
\pdef{Version} & Specifies a version, semantic versioning \texttt{major.minor.patch}. \\ \hline
\end{tabularx}

\newpage

\bibliographystyle{IEEEtran}
\bibliography{bibliography}

\newpage

\section{Revision History}
\subsection{Amendments}

\noindent\begin{tabularx}{\textwidth}{| p{1cm} | p{3cm} | p{2cm} | X | p{4cm} |} \hline
\rowcolor{gray!33} No. & Date & Version & Subject of Amendments & Author \\ \hline

1 & 2026-10-22 & 5.2.2 & Document created from the Arrowhead LaTeX templates for the 5.2.2 Java Spring implementation, including the local extensions of this workspace. & Jerker Delsing \\ \hline
2 & 2026-10-30 & 5.2.2 & Checked against the primary Java source tree (Eclipse-Arrowhead-java): services, operations, paths and data models unchanged. & Claude Opus 5.5 (draft) \\ \hline
3 & 2026-11-09 & 5.3.0 & New documentation edition; primary-source references and Go/Rust implementation scope; 5.2.2 retained. & Documentation maintenance \\ \hline
\end{tabularx}

\subsection{Quality Assurance}

\noindent\begin{tabularx}{\textwidth}{| p{1cm} | p{3cm} | p{2cm} | X |} \hline
\rowcolor{gray!33} No. & Date & Version & Approved by \\ \hline

1 & & \arrowversion  &  \\ \hline

\end{tabularx}

\end{document}
